\chapter*{Tóm tắt}
Nội dung của đồ án hướng đến việc nghiên cứu, thiết kế và hiện thực một IoT gateway dạng mô-đun có khả năng khả cấu hình linh hoạt, nhằm giải quyết các hạn chế của các giải pháp gateway truyền thống trong các hệ thống IoT hiện nay. Trong bối cảnh mỗi kịch bản ứng dụng IoT đòi hỏi tập hợp giao thức và chuẩn giao tiếp khác nhau, việc phát triển các gateway tích hợp sẵn nhiều chức năng không cần thiết dẫn đến lãng phí tài nguyên và làm tăng chi phí triển khai.

Trong giai đoạn đồ án chuyên ngành, nhóm tập trung tìm hiểu các khái niệm nền tảng liên quan đến IoT gateway, làm rõ vai trò của gateway trong hạ tầng IoT, đồng thời phân tích các công nghệ, giao thức và chuẩn giao tiếp phổ biến đang được sử dụng. Trên cơ sở đó, một kiến trúc gateway định hướng mô-đun được đề xuất, trong đó phần cứng được thiết kế theo dạng lắp ghép nhằm cho phép lựa chọn các mô-đun phù hợp với từng kịch bản ứng dụng, trong khi phần mềm hỗ trợ quy trình cấu hình tối giản và thân thiện với người dùng.

Sang giai đoạn đồ án tốt nghiệp, nhóm tập trung hoàn thiện sản phẩm dựa trên các kết quả đã đạt được, bao gồm tích hợp các thành phần phần cứng lên mạch in và hoàn thiện hệ thống phần mềm (firmware, software). Từ phiên bản nguyên mẫu ban đầu, nhóm tiến hành phân tích ưu điểm và hạn chế còn tồn tại, từ đó hiệu chỉnh đặc tả kỹ thuật và xây dựng kế hoạch chi tiết cho quá trình hoàn thiện hệ thống.

Sau khi hoàn thiện, hệ thống được kiểm thử về chức năng và hiệu năng. Kết quả cho thấy sản phẩm vận hành ổn định trong các điều kiện thử nghiệm, đáp ứng các yêu cầu chức năng đã đề ra và đảm bảo khả năng giao tiếp giữa các mô-đun phần cứng. Hệ thống cũng thể hiện tính linh hoạt trong cấu hình và mở rộng, phù hợp với nhiều kịch bản ứng dụng IoT khác nhau. Những kết quả này xác nhận tính khả thi của giải pháp đề xuất, đồng thời tạo nền tảng cho việc tiếp tục phát triển thành sản phẩm có khả năng ứng dụng thực tế.